\chapter{总结与展望}
\label{chp:conclusion}

\section{论文总结}

图推理任务要求模型同时完成结构理解、约束满足、算法选择和结果表达等多个环节，对大语言模型的综合求解能力提出了较高要求。现有方法虽然已经在部分图任务上展现出一定效果，但仍普遍存在端到端推理链条过长、知识调用缺乏中间组织以及程序生成结果不稳定等问题，导致模型在复杂图任务、强约束任务和动态图任务中容易出现理解偏差、实现错误或格式失配。围绕上述问题，本文提出了一种规划驱动的多智能体图推理框架 PlanGraph，通过在问题理解与程序执行之间引入显式伪代码计划，将规划、检索、编码、验证和修复等环节组织为一个相互衔接的求解闭环，从而提升复杂图推理任务的可靠性、稳定性与可解释性。

本文围绕所提方法开展了较为系统的研究，主要工作和研究结论如下。

（1）针对大语言模型在图推理任务中的主要困难进行了系统分析。本文从任务特性出发，梳理了图推理问题在结构表示、约束建模、算法映射和程序执行等方面的特殊性，指出图任务并不仅仅是一般文本推理的简单延伸，而是一类同时依赖结构化理解和可执行求解的复杂任务。进一步地，本文总结出当前大语言模型在该类任务中面临的几个关键瓶颈，包括端到端直接生成时推理负担过重、外部知识调用缺乏结构化引导、程序生成与调试过程稳定性不足等。这一分析为后续方法设计提供了明确的问题导向，也说明了仅依赖更大参数规模并不足以从根本上解决图推理中的可靠性问题。

（2）设计并实现了以规划为核心的多智能体图推理框架 PlanGraph。本文围绕图推理任务的求解流程，构建了问题智能体、规划智能体、检索智能体、编码智能体和推理智能体等功能角色，使系统能够按照“理解问题—生成计划—调用知识—生成程序—执行验证—反馈推理”的顺序逐步完成求解。与传统的端到端答案生成方式相比，该框架将原本隐含在模型内部的一次性推理过程显式拆解为若干相互协同的中间阶段，使各模块具有更清晰的职责边界，也使系统整体更便于分析、调试和扩展。实验与案例分析表明，这种规划驱动的组织方式能够有效降低复杂图任务中的错误传播风险。

（3）提出了基于伪代码计划的知识检索机制。为了缓解大语言模型在图算法调用和 API 选择上的不稳定性，本文没有直接使用原始问题文本进行知识检索，而是利用规划阶段生成的伪代码步骤作为中间查询表示，使检索内容更贴近算法流程与实现意图。与直接检索相比，这种方式能够更准确地定位图构建、路径搜索、匹配求解、流网络处理以及约束检查等相关知识，从而为编码阶段提供更可靠的实现依据。第五章的机制分析表明，规划驱动检索在知识命中率、覆盖率以及首次验证通过率等指标上均优于其他查询方式，说明中间规划不仅改善了推理过程，也实质性提升了知识调用质量。

（4）构建了程序执行验证与反馈推理闭环。针对图推理任务中“程序表面上能够运行但结果并不满足任务要求”的常见问题，本文在程序生成之后引入执行验证环节，通过测试样例、约束检查和格式校验对候选程序进行自动评估；当发现错误时，再将异常信息反馈给推理智能体进行局部修正或重新生成。该机制使系统不仅能够输出程序，还能够对程序是否真正满足任务目标进行二次确认。实验结果表明，验证与推理模块对整体性能提升具有稳定贡献，尤其在路径搜索、组合优化和强约束图任务中贡献更为显著。这表明，对于图推理这类可程序化求解的问题，执行反馈比单纯依赖自然语言自我反思更能提供可信的纠错依据。

（5）在多个公开图推理 benchmark 上系统验证了所提方法的有效性。本文在 GraphArena、NLGraph、GraphWiz、LLM4DyG 和 GraphInstruct 等数据集上对 PlanGraph 进行了实验评估，并从总体结果、稳定性、任务类别、复杂度分层、参数敏感性、工程开销、鲁棒性、消融实验和案例分析等多个角度展开了较为全面的验证。实验结果表明，PlanGraph 在多数 benchmark 上取得了最优或接近最优的性能，尤其在组合优化、路径搜索和结构约束较强的任务中表现出更为稳定的优势；同时，规划、检索与验证三者之间形成的协同闭环，是推动性能提升和稳定性增强的关键来源。相关分析也表明，PlanGraph 在动态图场景中仍存在进一步优化空间，这为后续研究提供了改进方向。

综上，本文研究表明：大语言模型在图推理任务中实现稳定可靠求解的关键，并不在于记忆更多孤立的图算法结论，而在于构建符合任务本质的系统化求解流程。本文提出的 PlanGraph 将规划、知识调用与执行验证组织为较为完整的求解链路，为大语言模型处理图结构任务提供了一种具有可解释性、可扩展性与工程可实现性的解决方案。

\section{论文展望}

尽管本文围绕规划驱动的多智能体图推理进行了较为系统的研究，并取得了较好的实验结果，但从更广泛的图智能应用与大模型系统发展角度看，当前工作仍然存在进一步扩展和深化的空间。未来可以从以下几个方面继续开展研究。

（1）面向动态图与异构图任务进一步增强规划表示能力。当前 PlanGraph 的规划模块已经能够较好支持经典图算法和部分动态图任务，但对复杂时间依赖、动态图演化过程以及异构节点和边类型之间关系的表达仍然相对粗粒度。未来可考虑设计更细粒度的结构化规划语言，使规划结果不仅包含算法步骤，还能够显式表示时间顺序、约束来源、状态变化和多类型关系之间的交互逻辑，从而增强系统对复杂图场景的适配能力。

（2）构建更丰富的外部知识库与更强的检索选择机制。本文当前所使用的知识主要来源于图计算文档和相关 API 说明，虽然已经能够支持多数实验任务，但面对更复杂、更开放的真实问题时，知识来源仍显有限。后续可以进一步引入算法教材、经典代码仓库、高质量技术问答以及专用图推理工具文档，并结合重排序、工具选择和多路检索融合等策略，提升系统在不同任务下调用外部知识的准确性与鲁棒性。

（3）研究更系统的程序验证与可信性保障方法。本文当前采用的小样例执行验证与错误修复机制能够有效发现部分实现问题，但在边界条件覆盖、约束完备性和形式化正确性方面仍有提升空间。未来可以将自动测试生成、约束求解器、静态分析方法以及形式化验证技术引入图推理系统，使程序验证不再局限于经验性样例检查，而是向更系统、更严格的正确性保障方向发展，以满足高可靠应用场景的需求。

（4）推动方法向真实图应用场景迁移。本文的实验主要建立在公开 benchmark 基础之上，这能够验证方法本身的有效性，但距离真实应用仍存在一定差距。未来可进一步将 PlanGraph 拓展到知识图谱问答、软件依赖分析、社交关系推断、网络运维诊断以及复杂流程优化等实际任务中，考察系统在开放输入、长链约束和多工具协同条件下的表现，并据此完善框架的工程化能力。

（5）探索多模型、多工具协同的混合推理系统。随着大语言模型、图神经网络、外部求解器和程序分析工具的不断发展，未来图推理系统很可能不再依赖单一模型完成全部任务，而是通过不同模型与工具的协同形成更高效的混合推理机制。后续研究可以考虑让语言模型负责任务理解与规划，让专用图模型负责结构表征，让外部求解器承担精确计算与验证，从而构建性能更强、效率更高且更具可控性的图推理系统。

综上所述，规划驱动的多智能体图推理在图任务上展现出较为广阔的研究前景。随着大语言模型能力持续增强以及外部工具生态不断完善，围绕结构化规划、知识调用、程序执行和结果验证展开的融合研究，有望进一步推动图推理系统朝着更可靠、更通用和更具工程实用价值的方向发展。
